
\begin{revision}
\begin{figure*}[t]
\centering
\begin{tikzpicture}[
  font=\footnotesize,
  stage/.style={draw,rounded corners=2pt,align=center,text width=2.6cm,
    minimum height=15mm,fill=blue!5},
  flow/.style={-Latex,semithick},
  node distance=7mm
]
\node[stage] (observe) {1. Observe\\telemetry + scenario};
\node[stage,right=of observe] (reason) {2. Propose\\retrieval + assessors};
\node[stage,right=of reason] (govern) {3. Govern\\policy ladder};
\node[stage,right=of govern] (bind) {4. Bind\\decision + evidence hash};
\node[stage,right=of bind] (audit) {5. Audit\\HUMAN gate + anchor};
\draw[flow] (observe) -- (reason);
\draw[flow] (reason) -- (govern);
\draw[flow] (govern) -- (bind);
\draw[flow] (bind) -- (audit);
\draw[flow,dashed] (audit.south) to[bend left=25]
  node[below]{receipt, decision, and evidence identifiers} (observe.south);
\end{tikzpicture}
\caption{\rev{Five-stage execution schematic. Model output remains untrusted until
the external policy ladder admits it. The complete exported workflow and its
node-level configuration are provided in the supplement.}}
\label{fig:five-stage-pipeline}
\end{figure*}
\end{revision}
